\documentclass{oci}
\usepackage[utf8]{inputenc}
\usepackage{lipsum}
\usepackage{tikz}
\usepackage{subcaption}

\title{Ejemplo}

\newcommand{\megawoman}{\scalebox{0.065}{\includegraphics{megawoman}}}
\newcommand{\robot}{\scalebox{0.1}{\includegraphics{robot}}}

\begin{document}
\begin{problemDescription}
La reconocida empresa de videojuegos Jockeycom está desarrollando una nueva
entrega de su aclamada franquicia Megawoman.
En esta nueva versión, nuestra heroína Mega Woman debe atravesar un escenario
plagado de robots bombarderos y así poder salvar a su amigo Uno.

El escenario está compuesto de $N$ casillas numeradas de 0 a $N-1$ que contiene
$R$ robots numerados de 0 a $R-1$.
Mega Woman corre siempre hacia la derecha con una velocidad constante de $v$
casillas por segundo, mientras que cada uno de los robots vuela por el
escenario en una sola dirección, izquierda o derecha.
El robot número $i$ vuela a una velocidad constante de $u_i$ casillas por
segundo.
El objetivo de Ultra Woman es escapar del escenario lográndolo cuando su
posición está más a la derecha que la última casilla.

Al comienzo de cada segundo, exceptuando el momento inicial del juego, los
robots lanzan una bomba en la casilla en la que están.
Si Ultra Woman se encuentra en ese momento dentro de la misma casilla será
alcanzada por la bomba y perderá el juego.
Los robots solo lanzan una bomba al inicio de cada segundo y por lo tanto Ultra
Woman puede pasar por debajo de ellos sin recibir daño si lo hace en el momento
indicado.
Por ejemplo, suponiendo que la velocidad de Ultra Woman es 3 casillas por
segundo, las siguientes figuras muestran como Ultra Woman puede pasar por debajo
de un robot con velocidad 2 sin ser dañada.

\begin{figure}[h]
\begin{subfigure}[c]{0.45\textwidth}
\begin{tikzpicture}[scale=1.2]
\node at (1.5, 0.2) {\megawoman};
\node at (3.5, 0.75) {\robot};
\draw[step=1.0,gray,thin] (0.01,0) grid (4.99,1);
\end{tikzpicture}
\caption{Momento inicial}
\end{subfigure}
~
\begin{subfigure}[c]{0.45\textwidth}
\begin{tikzpicture}[scale=1.2]
\node at (4.5, 0.2) {\megawoman};
\node at (1.5, 0.75) {\robot};
\draw[step=1.0,gray,thin] (0.01,0) grid (4.99,1);
\end{tikzpicture}
\caption{Situación 1 segundo después}
\end{subfigure}
\end{figure}

Dada la posición inicial de Ultra Woman y la de cada uno de los robots, a los
desarrolladores de Jockeycom les gustaría determinar si es posible para Ultra
Woman escapar del escenario.
?`Podrías ayudarlos?
\end{problemDescription}

\begin{inputDescription}
  La primera línea de la entrada contiene un entero $N$ ($0<N\leq 10^5$),
  correspondiente a la cantidad de casillas en el escenario.
  La siguiente línea contiene dos enteros $w$ y $v$ ($0\leq w< N, 0<v\leq 10^9$)
  correspondientes la casilla donde se encuentra inicialmente Ultra Woman y a la
  velocidad con la que se mueve.
  Posteriormente sigue una línea con un entero $R$ indicando la cantidad de
  robots en el escenario.
  A continuación siguen $R$ filas describiendo cada uno de los robots.
  Cada fila contiene dos enteros $r_i$ y $v_i$ ($0\leq 0 < N, -10^9\leq v_i\leq
  10^9$), correspondientes respectivamente a la posición inicial y la velocidad
  del robot.
  La velocidad será positiva si el robot se mueve hacia la derecha y negativa si
  se mueve hacia la izquierda.
\end{inputDescription}

\begin{outputDescription}
La salida debe corresponder a una única línea conteniendo la palabra, $SI$ o
$NO$, dependiendo de si Ultra Woman puede escapar del escenario sin ser
alcanzada por una bomba.
\end{outputDescription}

\begin{scoreDescription}
  \score{20} Se probarán varios casos donde la velocidad de Ultra Woman es 1,
  todos los robots se mueven hacia la derecha con velocidad 2 y el número de
  robots es menor o igual que $10^5$ ($u_i=2, v=1, 0\leq R\leq 10^5$).
  \score{20} Se probarán varios casos donde la velocidad de Ultra Woman es 1,
  todos los robots se mueven hacia la izquierda con velocidad 1 y la cantidad de
  robots es menor que $10^5$ ($u_i=1, v=1, 0\leq R\leq 10^5$).
  \score{30} Se probarán varios casos en que el número de robots es menor que 50
  y sin restricciones sobre las velocidades ($0\leq R\leq 50$).
  \score{30} Se probarán varios donde el número de robots es menor o igual que
  $10^5$ y sin restricciones adicionales ($0\leq R\leq 10^5$).
\end{scoreDescription}

\begin{sampleDescription}
\sampleIO{sample-1}
\sampleIO{sample-2}
\end{sampleDescription}

\end{document}